% =====================================================================
%  IEEE-STYLE CONFERENCE REPORT  (self-contained)
%  MLE-AI-201 Machine Learning -- Final Project
%
%  Figures:  python -m report.make_figures      (from the repo root)
%  Build:    tectonic report.tex                 (from report/)
%     or:    pdflatex report.tex ; pdflatex report.tex
% =====================================================================
\documentclass[10pt,twocolumn,a4paper]{article}

% ---- Geometry: IEEE-like margins / two narrow columns ----------------
\usepackage[a4paper,top=0.75in,bottom=1.0in,left=0.68in,right=0.68in,columnsep=0.25in]{geometry}

% ---- Times-like font (standard, ships with every TeX distro) ---------
% XeTeX/LuaTeX (tectonic): a Unicode Times clone covers Azerbaijani letters.
% pdfLaTeX: Times via mathptmx; the schwa is drawn as a rotated "e".
\usepackage{iftex}
\usepackage{graphicx}
\iftutex
  \usepackage{mathptmx}
  \usepackage{fontspec}
  \setmainfont{texgyretermes}[Extension=.otf,UprightFont=*-regular,
    BoldFont=*-bold,ItalicFont=*-italic,BoldItalicFont=*-bolditalic]
\else
  \usepackage{mathptmx}
  \usepackage[T1]{fontenc}
  \usepackage[utf8]{inputenc}
  \DeclareUnicodeCharacter{0259}{\rotatebox[origin=c]{180}{e}}
  \DeclareUnicodeCharacter{018F}{\rotatebox[origin=c]{180}{E}}
\fi

% ---- Common, dependency-free packages --------------------------------
\usepackage{graphicx}
\usepackage{amsmath,amssymb}
\usepackage{booktabs}
\usepackage{array}
\usepackage{titlesec}
\usepackage{caption}
\usepackage{cite}
\usepackage[hidelinks]{hyperref}
\usepackage{url}
\usepackage{enumitem}
\setlist{nosep,leftmargin=1.4em}

\graphicspath{{figures/}}

% ---- IEEE-style section headings -------------------------------------
\renewcommand{\thesection}{\Roman{section}}
\renewcommand{\thesubsection}{\Alph{subsection}}
\titleformat{\section}[block]{\centering\normalsize\scshape}{\thesection.}{0.5em}{}
\titleformat{\subsection}[block]{\normalsize\itshape}{\thesubsection.}{0.5em}{}
\titlespacing*{\section}{0pt}{1.4ex plus 1ex minus .2ex}{1ex}
\titlespacing*{\subsection}{0pt}{1.0ex plus 1ex minus .2ex}{0.6ex}

% ---- Abstract / Index-Terms in the IEEE boldface-dash style ----------
\newcommand{\abstractkw}[1]{\noindent\textbf{\textit{#1}}}

% ---- Team member whose surname is still to be filled in ---------------
\newcommand{\DavudSurname}{[Surname]}

\captionsetup{font=small,labelsep=period}
\setlength{\parindent}{1em}

% =====================================================================
\begin{document}

% ------------------------- TITLE BLOCK --------------------------------
\twocolumn[
  \begin{@twocolumnfalse}
  \begin{center}
    {\LARGE\bfseries Decision Trees and Pegasos SVMs from Scratch:\\Price and Price-Tier Prediction on bina.az Listings\par}
    \vspace{1.2em}
    % Authors: Surname, Given name --- alphabetical by surname.
    {\large
      Bayramov, Bailar\quad \DavudSurname, Davud\quad Hajiyev, Sanan\quad
      Khamiyev, Sabuhi\quad Samedov, Farid\par}
    \vspace{0.5em}
    {\itshape MLE-AI-201 Machine Learning, AI Academy --- Cohort II\par}
    \vspace{1.4em}
  \end{center}
  \begin{quote}
    \small
    \abstractkw{Abstract---}%
    We predict residential sale prices in Baku from 100,775 raw bina.az
    listings and frame the same data as two tasks: price regression (Task~A)
    and a binary premium/standard price tier split at the training median of
    209,000~AZN (Task~B). After URL deduplication, residential/Baku filtering,
    parsing of Azerbaijani text fields and fixed-rule outlier removal, 57,363
    listings with 65 encoded features remain; the leakage columns
    \texttt{unit\_price} and \texttt{total\_price} are removed. We implement a
    CART-style decision tree (Gini, entropy, MSE) and a soft-margin SVM trained
    with the Pegasos sub-gradient method, with random Fourier features (RFF) as
    the non-linear extension, entirely in NumPy, and benchmark them against
    scikit-learn with matched settings. On the held-out test split our tree
    reaches RMSE 109,164~AZN ($R^2=0.760$) versus 119,376~AZN for the
    independently tuned sklearn tree, and F1 0.929 for the price tier (sklearn
    0.929). Our linear Pegasos SVM reaches F1 0.881 versus 0.903 for linear
    SVC while training $\sim$100$\times$ faster; RFF does not beat the linear
    model. Five-fold cross-validation confirms the ranking. Most squared error
    comes from a small set of luxury listings that the available columns cannot
    explain. Code:
    \href{https://github.com/EmbeddedDesignEngineer/ML-Final}{\texttt{github.com/EmbeddedDesignEngineer/ML-Final}}
    (tag \texttt{v1.0-final}).

    \vspace{0.6em}
    \abstractkw{Index Terms---}decision tree, support-vector machine,
    Pegasos, random Fourier features, regression, classification,
    real-estate price prediction.
  \end{quote}
  \vspace{1.0em}
  \end{@twocolumnfalse}
]

% =====================================================================
\section{Introduction}
Online listing platforms such as bina.az, Azerbaijan's largest real-estate
site, expose asking prices for tens of thousands of properties. A model that
estimates a fair price from size, location and condition helps buyers and
sellers spot mispriced listings. The data are a raw scrape: Azerbaijani
column names, numbers stored as text (``145~m\textsuperscript{2}'',
``7 / 9''), duplicates, impossible values and columns that leak the price.

Our contributions are: (i) a documented, leakage-free cleaning and encoding
pipeline whose every statistic is fitted on training rows only;
(ii) a NumPy decision tree supporting Gini, entropy and MSE with four working
stopping controls; (iii) a NumPy Pegasos SVM with margin and support-vector
diagnostics and an RFF kernel approximation; (iv) a matched benchmark against
scikit-learn with five-fold uncertainty, timing and subgroup error analysis;
and (v) two bonuses: a NumPy Ridge baseline and a bootstrap ensemble of our own
trees. One command, \texttt{python -m src.run\_all}, reproduces every
headline number.

\section{Background}
A classification tree chooses, at each node, the feature/threshold pair that
maximises the impurity decrease
$\Delta I = I(\text{parent}) - \sum_{c} \tfrac{N_c}{N} I(c)$, with Gini
$G = 1-\sum_k p_k^2$ or entropy $H=-\sum_k p_k\log_2 p_k$; a regression tree
uses the node variance $\tfrac{1}{N}\sum_i (y_i-\bar y)^2$~\cite{cart}.
The soft-margin linear SVM minimises the regularised hinge loss
\begin{equation}
  J(\mathbf{w},b)=\tfrac{\lambda}{2}\lVert\mathbf{w}\rVert^{2}
  +\tfrac{1}{n}\sum_{i=1}^{n}\max\!\big(0,\,1-y_i(\mathbf{w}^{\top}\mathbf{x}_i+b)\big),
\end{equation}
which Pegasos~\cite{pegasos} solves without a QP solver by stochastic
sub-gradient steps. Random Fourier features~\cite{rff} approximate the RBF
kernel $k(\mathbf{x},\mathbf{x}')=\exp(-\gamma\lVert\mathbf{x}-\mathbf{x}'\rVert^2)$
with an explicit map, so the linear solver can learn a non-linear boundary.

\section{Data and Preprocessing}
\subsection{Cleaning}
The Kaggle dump (\texttt{house\_sale.csv}, 100,775 rows, 51 columns,
SHA-256 \texttt{b6a4c67e\ldots}) is cleaned by \texttt{src/data\_prep.py}.
Each rule is fixed in advance (no data quantiles), so no statistic can leak
from validation or test rows:
\begin{enumerate}
  \item Drop repeated scrapes of the same listing URL, keeping the latest
        (100,775$\to$64,454).
  \item Keep the residential categories \emph{Yeni tikili} (new build),
        \emph{Köhnə tikili} (old build) and \emph{Həyət evi/Bağ evi} (house)
        (58,452), and Baku only (58,257; 99.7\% of residential rows).
  \item Parse \emph{Sahə} (area), \emph{Mərtəbə} (``floor / total''),
        \emph{Torpaq sahəsi} (land, sot); map \emph{Təmir}, \emph{Çıxarış},
        \emph{İpoteka} (repair, bill of sale, mortgage) to 0/1. All prices are
        in AZN; any other currency raises an error.
  \item Structural fills: houses get floor 0 of 1, apartments land 0. Rows
        still missing an essential field are dropped (57,587).
  \item Domain outlier rules: price 10k--5M~AZN, area 15--1,000~m\textsuperscript{2},
        price/m\textsuperscript{2} 200--15,000 (filter only, never a feature),
        1--15 rooms, land $\le$100 sot, $0\le$ floor $\le$ total floors,
        coordinates inside Azerbaijan (57,363 rows).
\end{enumerate}

\textbf{Leakage.} \texttt{unit\_price} (price per m\textsuperscript{2}) and
\texttt{total\_price} (equal to \texttt{price} in 99.99\% of rows) are
arithmetic functions of the target and are dropped. We also drop post-listing
signals (\texttt{views}, \texttt{updated}), identifiers and personal data
(owner/shop names, URLs, raw address), and free text.

\subsection{Features, tasks and split}
The 17 cleaned columns become 65 features: 14 numeric/binary inputs,
distance to Fountain Square (haversine) and floor ratio, a one-hot category
(3), and one-hot districts with at least 200 \emph{training} listings (45)
plus \texttt{loc\_other} for rare or unseen districts. The encoder is fitted on
the training rows only. Rows are split 70/15/15 with seed 42, stratified on
price deciles (Table~\ref{tab:data}). The tier label is
$\mathbf{1}[\text{price} > 209{,}000]$, where 209,000~AZN is the training
median; validation/test use the same threshold. SVM inputs are standardised
with training means and standard deviations; trees use raw features.

\begin{table}[t]
  \centering
  \caption{Dataset summary after cleaning.}
  \label{tab:data}
  \begin{tabular}{@{}lr@{}}
    \toprule
    Property & Value \\
    \midrule
    Raw rows / after URL deduplication  & 100,775 / 64,454 \\
    Rows (after cleaning)       & 57,363 \\
    Features used               & 65 \\
    Train / Val / Test split    & 40,155 / 8,604 / 8,604 \\
    Tier threshold (train median) & 209,000 AZN \\
    Premium share train / val / test (\%) & 49.71 / 49.73 / 49.72 \\
    \bottomrule
  \end{tabular}
\end{table}

\subsection{Exploration}
Prices are heavily right-skewed (Fig.~\ref{fig:eda}); $\log(\text{price})$
is close to symmetric. Area is the dominant driver, and location matters:
coastal and central districts (Bilgəh, Badamdar, Səbail) command far higher
prices per m\textsuperscript{2}. By construction the tier label is balanced
(49.7\% premium).

\begin{figure}[t]
  \centering
  \includegraphics[width=\linewidth]{price_distribution.png}
  \caption{Cleaned price distribution in AZN (left) and natural log of price
  (right).}
  \label{fig:eda}
\end{figure}

% =====================================================================
\section{Method}
\subsection{Decision tree (from scratch)}
\texttt{DecisionTree} in \texttt{src/decision\_tree.py} mirrors sklearn's
API (\texttt{fit}, \texttt{predict}, \texttt{predict\_proba}) for both tasks,
with the criterion and four stopping controls as constructor arguments.
At each node every feature is sorted once and
thresholds are placed midway between distinct adjacent values; a row goes left
if $x_j\le t$. Prefix class counts (classification) or centred prefix sums and
squares (regression) evaluate all thresholds in $O(n)$ after the
$O(n\log n)$ sort. The split with the largest gain wins; exact ties keep the
first feature and lowest threshold. Growth stops at a pure node, at
\texttt{max\_depth}, when a node has fewer than \texttt{min\_samples\_split}
rows, when a child would have fewer than \texttt{min\_samples\_leaf}, or when
the weighted gain $\tfrac{N_t}{N}\Delta I$ is below
\texttt{min\_impurity\_decrease} (sklearn's definition). Leaves predict the
majority class or the mean; \texttt{predict\_proba} returns leaf class
frequencies. Feature importance sums the weighted gains per feature and
normalises them. Unit tests check a tiny, fully determined example against
sklearn's thresholds, impurities and predictions, and that every stopping
control changes the tree.

\subsection{Support-vector machine (from scratch)}
\texttt{src/svm.py} implements \texttt{PegasosSVM}. Labels $\{0,1\}$ are mapped
to $\{-1,+1\}$; weights start at zero. At step $t$ a random row $i$ is drawn
and, with $\eta_t=1/(\lambda t)$ (Pegasos) or a constant $\eta$,
\begin{equation}
  \mathbf{w}\leftarrow(1-\eta_t\lambda)\mathbf{w}+
  \eta_t\,y_i\mathbf{x}_i\,\mathbf{1}[y_i(\mathbf{w}^\top\mathbf{x}_i+b)<1],
\end{equation}
with $b\leftarrow b+\eta_t y_i$ on margin violations (the bias is not
regularised), followed by the Pegasos projection onto
$\lVert\mathbf{w}\rVert\le1/\sqrt\lambda$. The full objective $J$ is logged
at 51 checkpoints. After training we report $\mathbf{w}$, the margin width
$2/\lVert\mathbf{w}\rVert$ and the margin-active training points
($y_if(\mathbf{x}_i)\le1$, the support vectors of the primal). Features are
standardised first because the hinge loss is scale-sensitive.

\textbf{Non-linearity.} We chose random Fourier features (option ii):
$z(\mathbf{x})=\sqrt{2/D}\cos(\mathbf{x}W+\mathbf{c})$,
$W_{jk}\sim\mathcal N(0,2\gamma)$, $c_k\sim U(0,2\pi)$. Kernelised Pegasos
would store a coefficient per training row and cost $O(n)$ kernel
evaluations per update and per prediction; with 40k rows RFF keeps the
cheap linear update and a fixed memory budget $D$.

\subsection{Bonus models}
\textbf{Ridge} (\texttt{src/linear.py}) solves
$(X_c^\top X_c+\alpha I)\mathbf{w}=X_c^\top y_c$ on centred data via an SVD,
leaving the intercept unpenalised. \textbf{Bagging}
(\texttt{src/ensemble.py}) averages our own trees, each fitted on a bootstrap
sample with optional feature subsampling.

% =====================================================================
\section{Experimental Setup}
\textbf{Metrics.} Task A: RMSE, MAE and $R^2$ in AZN. Task B: precision,
recall, F1 for the premium class, ROC-AUC from decision scores or leaf
probabilities, and the confusion matrix. Metrics are our own NumPy functions
in \texttt{src/evaluate.py}, tested against sklearn.

\textbf{Selection protocol.} Every hyperparameter is chosen on the
validation split; the test split is scored once with frozen settings.
Trees: depth $\{2,4,8,12\}\times$ leaf $\{1,5,20\}$, plus Gini vs entropy for
Task B, chosen by validation RMSE / F1. SVM: 33 predeclared configurations
(Table~\ref{tab:svmgrid}), chosen by validation F1 (ROC-AUC tie-break).
Ridge: $\alpha\in\{0.1,1,10,100,1000\}$. The custom and sklearn models are
selected independently from the same grids.

\begin{table}[t]
  \centering\small
  \caption{SVM search grid (all fits on the full training split).}
  \label{tab:svmgrid}
  \begin{tabular}{@{}p{0.23\linewidth}p{0.6\linewidth}r@{}}
    \toprule
    Family & Grid & Fits \\
    \midrule
    Ours, linear & $\lambda\in\{10^{-4},10^{-3},10^{-2}\}$ $\times$ $T\in\{5\text{k},20\text{k}\}$ $\times$ \{Pegasos, $\eta{=}0.01$\} & 12\\
    Ours, RFF & same $\lambda$ $\times$ $\gamma\in\{0.01,0.1\}$ $\times$ $D\in\{128,512\}$, $T{=}20$k & 12 \\
    SVC linear & same $\lambda$, $C=1/(n\lambda)$ & 3 \\
    SVC RBF & same $\lambda\times\gamma$ & 6 \\
    \bottomrule
  \end{tabular}
\end{table}

\textbf{Baselines.} \texttt{DecisionTreeRegressor/Classifier} with the same
criterion, depth and leaf size; \texttt{SVC(kernel="linear")} and
\texttt{SVC(kernel="rbf")} inside a \texttt{Pipeline} with
\texttt{StandardScaler}. Our loss averages the hinge term while SVC sums it,
so $C=1/(n_\text{train}\lambda)$ matches regularisation strength.

\textbf{Uncertainty.} Five-fold CV on the train+validation pool (test rows
excluded) refits the encoder, scaler and tier threshold inside every fold
and reports mean $\pm$ sample std (ddof=1) for the frozen settings.
Seeds are fixed at 42. Timings are single-run wall-clock measurements on a
laptop CPU.

% =====================================================================
\section{Results}
\subsection{Tuning the tree}
Fig.~\ref{fig:tree} shows the bias--variance story. Training error keeps
falling with depth, while validation RMSE flattens after depth 8 and, with
leaf size 1, \emph{rises} again (134k at depth 12): the tree memorises single
expensive listings. A leaf size of 5 regularises enough for depth 12 to win
(validation RMSE 124,741). For Task B, Gini beats entropy at every depth $\ge4$
with leaf size 1 (0.9285 vs 0.9223 at depth 12); the F1 gap between training and
validation widens with depth but validation F1 still improves, so depth 12
with Gini is selected.

\begin{figure}[t]
  \centering
  \includegraphics[width=\linewidth]{tree_curves.png}
  \caption{Validation curves of our tree. Left: RMSE vs depth for three leaf
  sizes. Right: F1 vs depth, Gini vs entropy. Dashed lines are training
  scores.}
  \label{fig:tree}
\end{figure}

\subsection{Tuning the SVM}
With $\lambda=10^{-4}$ the constant schedule reduces the objective from 1.0
to 0.309 within 20k updates, while the $1/(\lambda t)$ schedule starts with
huge steps ($\eta_1=10^4$), overshoots and is still at 2.69 after 20k
updates (Fig.~\ref{fig:svm}, left): 20k updates are only half an epoch, so
convergence is not established for that schedule. Increasing $\lambda$ widens
the margin from 0.45 to 0.71 (constant schedule) but lowers validation F1 from
0.887 to 0.871 (Fig.~\ref{fig:svm}, right): a wider margin is not
automatically a better classifier here. The selected linear model
($\lambda=10^{-4}$, $\eta=0.01$, 20k updates) has $\lVert\mathbf{w}\rVert=4.42$,
margin width 0.453 and 10,532 margin-active training points (26.2\%). For RFF,
$D=512$ beats $D=128$ and $\gamma=0.01$ beats $\gamma=0.1$ in every pairing;
the winner ($\lambda=10^{-3}$, $\gamma=0.01$, $D=512$) has 53\% margin-active
points, a sign of an under-optimised, high-dimensional model.

\begin{figure}[t]
  \centering
  \includegraphics[width=0.92\linewidth]{svm_objective.png}\\[0.3em]
  \includegraphics[width=0.92\linewidth]{svm_margin.png}
  \caption{Top: training objective of our linear SVM ($\lambda=10^{-4}$) for
  both schedules. Bottom: margin width $2/\lVert\mathbf{w}\rVert$ vs $\lambda$.}
  \label{fig:svm}
\end{figure}

\subsection{Benchmark against scikit-learn}
Table~\ref{tab:results} is the required results table; Fig.~\ref{fig:results}
visualises it. At \emph{matched} settings the two trees agree closely: all 18
configurations with depth $\le$4 give identical validation scores, and the
largest Task-B gap is 0.001 F1. Differences appear only in deep trees, where
equal-gain ties (sklearn permutes features at random, we keep the first) and
sklearn's float32 thresholds send a few rows down different branches; deep
regression trees amplify these into different leaves (custom vs sklearn at
depth 12, leaf 5: 124,741 vs 125,668 validation RMSE). The sklearn search
selected depth 8, so its test RMSE (119,376) is not a matched comparison;
our depth-12 tree is 8.6\% better on test. Ridge agrees with sklearn to
numerical precision, which validates that implementation.

\begin{table*}[t]
  \centering
  \caption{Headline results. ``Test'' is the held-out split (8,604 rows,
  scored once). ``5-fold CV'' is mean $\pm$ std over folds of the
  train+validation pool with the same frozen settings. Fit time is a single
  training run on 40,155 rows.}
  \label{tab:results}
  \small
  \begin{tabular}{@{}llccccr@{}}
    \toprule
    \multicolumn{7}{@{}l}{\textbf{Task A --- price regression (AZN)}}\\
    Model & Setting & Test RMSE & Test MAE & Test $R^2$ & 5-fold CV RMSE & Fit (s)\\
    \midrule
    Ours: decision tree & MSE, depth 12, leaf 5 & \textbf{109,164} & \textbf{45,555} & \textbf{0.760} & \textbf{119,675 $\pm$ 2,954} & 3.28\\
    sklearn \texttt{DecisionTreeRegressor} & MSE, depth 8, leaf 5 & 119,376 & 52,230 & 0.712 & 123,508 $\pm$ 2,965 & 0.21\\
    Ours: Ridge (bonus) & $\alpha=1000$ & 128,441 & 57,859 & 0.667 & 136,831 $\pm$ 5,550 & 0.07\\
    sklearn \texttt{Ridge} & $\alpha=1000$ & 128,441 & 57,859 & 0.667 & 136,831 $\pm$ 5,550 & 0.06\\
    \midrule
    \multicolumn{7}{@{}l}{\textbf{Task B --- price tier (premium = 1)}}\\
    Model & Setting & Precision & Recall & F1 & 5-fold CV F1 & Fit (s)\\
    \midrule
    Ours: decision tree & Gini, depth 12, leaf 1 & 0.929 & 0.930 & \textbf{0.929} & 0.925 $\pm$ 0.003 & 2.84\\
    sklearn \texttt{DecisionTreeClassifier} & Gini, depth 12, leaf 1 & 0.929 & 0.929 & \textbf{0.929} & 0.925 $\pm$ 0.003 & 0.32\\
    Ours: linear SVM (Pegasos) & $\lambda{=}10^{-4}$, $\eta{=}0.01$, 20k & 0.874 & 0.887 & 0.881 & 0.885 $\pm$ 0.005 & 0.30\\
    Ours: RFF SVM & $\lambda{=}10^{-3}$, $\gamma{=}0.01$, $D{=}512$ & 0.766 & 0.945 & 0.846 & 0.836 $\pm$ 0.009 & 1.17\\
    sklearn \texttt{SVC} linear & $C=0.249$ & 0.905 & 0.900 & 0.903 & 0.906 $\pm$ 0.002 & 43.1\\
    sklearn \texttt{SVC} RBF & $C=0.249$, $\gamma=0.01$ & 0.894 & 0.892 & 0.893 & 0.893 $\pm$ 0.001 & 54.2\\
    \bottomrule
  \end{tabular}
  \\[0.4em]
  \small Test ROC-AUC: tree 0.949 (sklearn 0.949), linear SVM 0.949, RFF 0.934, SVC linear 0.963, SVC RBF 0.964.
\end{table*}

\begin{figure}[t]
  \centering
  \includegraphics[width=\linewidth]{results_bars.png}
  \caption{Test metrics (bars; ours in blue), five-fold CV means (diamonds)
  and CV standard deviation (error bars).}
  \label{fig:results}
\end{figure}

For the SVM, linear SVC is 2.2 F1 points better than our Pegasos model.
Both optimise the same objective, but SVC solves it exactly (libsvm SMO)
while our 20k stochastic updates are half an epoch. SVC pays for this with
a 43~s fit versus 0.3~s, and with prediction through 10,765 stored support
vectors (2.6~s vs 0.0015~s on the test set). Exact RBF SVC is also better than
our RFF approximation, which is limited both by $D=512$ random features and
by incomplete optimisation. Neither non-linear model beats its linear
counterpart in F1, so the extra complexity is not justified on this data.

\subsection{Improvement study (validation only)}
Because the test split had already been seen, follow-up improvements were
judged on validation only (\texttt{--quality-study}, \texttt{--price-study}).
Averaging the Pegasos iterate and training for 20 epochs raised validation F1
to 0.907 (linear) and 0.898 (RFF with $\lambda=10^{-5}$), closing the gap to
SVC and confirming that the earlier gap was optimisation, not the model. For
Task A, a bagged ensemble of 48 of our trees (depth 32, bootstrap 0.8)
trained on $\log(\text{price}/\text{area})$ reduced validation RMSE from
124,741 to 97,157~AZN and MAE from 48,345 to 33,968~AZN ($R^2$ 0.822), and
validation MAE falls monotonically as trees are added (1$\to$24 trees:
50.2k$\to$37.3k). These are validation numbers and are not mixed into the
test table.

% =====================================================================
\section{Error Analysis and Interpretation}
\textbf{What the models use.} The tree's impurity importance is dominated by
area (0.63) and distance to the centre (0.16), followed by longitude,
latitude, total floors and land size (Fig.~\ref{fig:feat}). The largest
standardised weights of the linear SVM tell the same story: area ($+3.63$),
house category ($-0.99$), distance to centre ($-0.85$) and rooms ($+0.83$);
suburban districts such as Masazır have negative weights. These are
conditional associations of correlated features, not causal effects.

\begin{figure}[t]
  \centering
  \includegraphics[width=\linewidth]{tree_reg_importance.png}
  \caption{Top impurity-decrease importances of our regression tree.}
  \label{fig:feat}
\end{figure}

\textbf{Luxury outliers.} Listings above the training 90th percentile
(455,000~AZN; 856 test rows) have a tree MAE of 177,315~AZN versus 30,998 for
the rest, with a negative bias ($-56,806$): expensive homes are
underpredicted. In the validation study the worst 1\% of rows produce
70--74\% of the squared error. Most are real villas in Bilgəh, Badamdar or
Səbail at 3--8k~AZN/m\textsuperscript{2}; nothing in the columns (sea view,
finish quality, exact address) separates them from ordinary houses. A few are
probable data errors (e.g.\ 600~m\textsuperscript{2} for 159k~AZN). This is
why RMSE stays high while median relative error is only 6--9\%.

\textbf{Sparse districts.} Districts with fewer than 200 training listings
are pooled into \texttt{loc\_other}; their MAE is 54,700 vs 44,673~AZN for
frequent districts, and their SVM error rate is 16\% vs 12\%.

\textbf{Classification errors} concentrate near the 209k threshold: only 7 of
856 expensive listings are misclassified by the tree. The RFF model trades
precision for recall (1,237 false positives vs 546 for the linear SVM).

% =====================================================================
\section{Discussion and Limitations}
\begin{itemize}
  \item The target is the \emph{asking} price, not a transaction price.
  \item The population is the filtered Baku residential market; conclusions
        do not cover commercial property or other cities.
  \item Five-fold CV quantifies variability of frozen settings, not nested
        model-selection uncertainty. The improvement study re-used the same
        validation split many times and is not an independent benchmark.
  \item Our Pegasos runs are short; the $1/(\lambda t)$ schedule did not
        converge at $\lambda=10^{-4}$.
  \item Deduplication uses listing URLs only; 446 exact feature duplicates
        remain (139 validation rows have a twin in training). Removing them
        changes validation MAE by under 1\%.
  \item Timings are single runs on one machine, not repeated benchmarks; our
        tree fits 9--15$\times$ slower than sklearn's Cython tree.
\end{itemize}

\section{Conclusion}
Our NumPy tree matches scikit-learn at matched settings and is the best
single model on both tasks (test RMSE 109k~AZN, F1 0.929). Our Pegasos SVM is
within 2 F1 points of exact SVC while training two orders of magnitude
faster, and longer averaged training closes most of that gap. The main
remaining error comes from luxury properties that the available features
cannot price. Next steps are richer location and quality features (exact
address, view, renovation level), deduplication by content, and a fresh
untouched test set for the improved ensemble.

\section*{Contributions}
\begin{itemize}[leftmargin=0.8em,label={}]\raggedright
  \item \textbf{Davud \DavudSurname}: data loading, cleaning, encoding and
        leakage audit; model-quality, price-error and diagnostics studies
        (\texttt{src/data\_prep.py}, \texttt{experiments/model\_quality.py},
        \texttt{price\_quality.py}, \texttt{price\_diagnostics/}).
  \item \textbf{Sanan Hajiyev}: decision tree, validation sweeps and sklearn
        parity tests (\texttt{src/decision\_tree.py},
        \texttt{experiments/tree\_experiments.py}).
  \item \textbf{Farid Samedov}: Pegasos SVM, RFF and the 33-configuration SVM
        study (\texttt{src/svm.py}, \texttt{experiments/svm\_*.py}).
  \item \textbf{Sabuhi Khamiyev}: metrics, evaluation, plots and subgroup
        analysis (\texttt{src/evaluate.py}, \texttt{src/evaluation\_plots.py}).
  \item \textbf{Bailar Bayramov}: Ridge bonus, shared runner and integration
        tests (\texttt{src/linear.py}, \texttt{src/run\_all.py},
        \texttt{experiments/pipeline.py}).
\end{itemize}
Details: \texttt{contribution\_report.md}.

\section*{Tools and Acknowledgements}
Data: Bina.az Sale Project dataset by S.~Memmedli on Kaggle~\cite{data}.
scikit-learn~\cite{sklearn} is used only for baselines. AI assistance: the
team used AI coding assistants (including Claude) for code review, test
suggestions, debugging, running the experiment pipeline and drafting and
formatting parts of this report and the slides. All numbers come from
executing the repository code; every member reviewed and is responsible for
their listed contribution.

% =====================================================================
\begin{thebibliography}{9}
\bibitem{pegasos}
S.~Shalev-Shwartz, Y.~Singer, N.~Srebro, and A.~Cotter,
``Pegasos: primal estimated sub-gradient solver for SVM,''
\emph{Mathematical Programming}, vol.~127, no.~1, pp.~3--30, 2011.
\bibitem{cart}
L.~Breiman, J.~Friedman, R.~Olshen, and C.~Stone,
\emph{Classification and Regression Trees}. Wadsworth, 1984.
\bibitem{rff}
A.~Rahimi and B.~Recht, ``Random features for large-scale kernel machines,''
in \emph{Proc. NeurIPS}, 2007, pp.~1177--1184.
\bibitem{sklearn}
F.~Pedregosa \emph{et al.}, ``Scikit-learn: Machine learning in Python,''
\emph{JMLR}, vol.~12, pp.~2825--2830, 2011.
\bibitem{data}
S.~Memmedli, ``Bina.az Sale Project,'' Kaggle dataset,
\url{https://www.kaggle.com/datasets/sehriyarmemmedli/binaaz-sale-project}.
\end{thebibliography}

\end{document}
